\section*{第 \ref{ch:g11:dna-replication} 章——DNA 复制}

\begin{solution}{exo:g11:dna-replication:1}
每条链都作为模板，据此合成一条新的互补链。
之所以能够如此，是因为 A 与 T 配对、G 与 C 配对：
一条链的序列决定了另一条链的序列，
因此单独一条链就携带了全部信息。
\end{solution}

\begin{solution}{exo:g11:dna-replication:2}
半保留复制：每个子代分子保留一条亲代链，
并新合成一条链。全保留复制：亲代分子保持完整，
旁边合成一个全新的双链分子。
\end{solution}

\begin{solution}{exo:g11:dna-replication:3}
\omterm{def:g10:cells-common-unit:cell}{细胞}周期的 S 期；每个\omterm{def:g10:cells-common-unit:organelle}{细胞核}中的 \omterm{def:g10:universal-dna:information}{DNA} 量加倍。
\end{solution}

\begin{solution}{exo:g11:dna-replication:4}
\omterm{def:g10:universal-dna:information}{DNA} 聚合酶逐个把与模板互补的\omterm{def:g10:universal-dna:nucleotide}{核苷酸}加到
正在延伸的链上。\omterm{prop:g11:dna-replication:forks}{复制叉}是移动的位点，
亲代链在此分开，
两条新链正在此处合成。
\end{solution}

\begin{solution}{exo:g11:dna-replication:5}
\texttt{CTAATGTCCG}，逐碱基对齐。
\end{solution}

\begin{solution}{exo:g11:dna-replication:6}
全保留复制，第 1 代：一条重带和一条轻带，没有杂合带
——但实际只观察到了杂合带。分散复制，第 2 代：
一条位于轻带与重带之间四分之一处的单带
——但实际观察到了两条不同的条带。每种模型都预测了
未曾观察到的条带。
\end{solution}

\begin{solution}{exo:g11:dna-replication:7}
$2^5 = 32$ 个分子，其中 2 个为杂合分子，30 个为轻分子：
轻带的强度是杂合带的 15 倍，且没有重带。
\end{solution}

\begin{solution}{exo:g11:dna-replication:8}
第 0 代：一条轻带。第 1 代：一条杂合带。
第 2 代：杂合带与重带强度相同——
恰好是原始实验的镜像。
\end{solution}

\begin{solution}{exo:g11:dna-replication:9}
$\num{2.5e8}/50 = \num{5e6}\,\unit{s}$，约 58 天——远长于
8 小时。\omterm{def:g10:universal-dna:information}{染色体}必须从许多复制起点同时开始复制，
每个起点打开两个\omterm{prop:g11:dna-replication:forks}{复制叉}。
\end{solution}

\begin{solution}{exo:g11:dna-replication:10}
两条链都需复制，每段 $\num{3.2e9}$ 个碱基对的 \omterm{def:g10:universal-dna:information}{DNA} 合计
为 $\num{6.4e9}$ 个碱基对，即 $\num{6.4e9}$ 个新\omterm{def:g10:universal-dna:nucleotide}{核苷酸}被连接（每复制一个碱基对，
连接一个\omterm{def:g10:universal-dna:nucleotide}{核苷酸}）。按每个\omterm{def:g10:universal-dna:nucleotide}{核苷酸} $10^{-9}$ 的错误率计算，约有 6 个错误。
\end{solution}

\begin{solution}{exo:g11:dna-replication:11}
打开一个复制泡时，起点两侧的模板都暴露出来；每个\omterm{prop:g11:dna-replication:forks}{复制叉}
都从起点向外复制，因此两叉方向相反，
共同覆盖整个区域。两个复制泡相遇时，它们的\omterm{prop:g11:dna-replication:forks}{复制叉}融合，
子代链首尾相接，
形成连续的分子。
\end{solution}

\begin{solution}{exo:g11:dna-replication:12}
细菌停止分裂（既不复制，也不分裂）。
伤口停止愈合，因为其修复需要\omterm{def:g10:cells-common-unit:cell}{细胞}分裂。神经\omterm{def:g10:cells-common-unit:cell}{细胞}不受影响，
它不再复制 \omterm{def:g10:universal-dna:information}{DNA}。一种能抑制快速分裂\omterm{def:g10:cells-common-unit:cell}{细胞}而不伤及不分裂\omterm{def:g10:cells-common-unit:cell}{细胞}的药物，
提示可用于抗击细菌感染，
或治疗分裂中的癌\omterm{def:g10:cells-common-unit:cell}{细胞}。
\end{solution}

\begin{solution}{exo:g11:dna-replication:13}
两条原始链从未被破坏：每一代中，它们各自再次与一条新的轻链配对，
因此杂合分子始终恰好有两个。
10 代之后：$2/2^{10} = 2/1024
\approx 0.2\%$——存在，但太淡，无法观察到。
\end{solution}

\begin{solution}{exo:g11:dna-replication:14}
$\num{6.4e9} \times 10^{-5} = \num{64000}$ 个错误每次分裂。若有两万个\omterm{def:g10:universal-dna:gene}{基因}，
则每次分裂都会有数百个\omterm{def:g10:universal-dna:gene}{基因}受损；
经过几次分裂后，\omterm{def:g10:cells-common-unit:cell}{细胞}便带有致死缺陷，
该\omterm{def:g10:cells-common-unit:cell}{细胞}系随之灭绝。
\end{solution}

\begin{solution}{exo:g11:dna-replication:15}
新一轮复制在上一轮尚未结束时即从起点开始，
\omterm{def:g10:universal-dna:information}{染色体}因此同时带有若干对嵌套的\omterm{prop:g11:dna-replication:forks}{复制叉}。
每次分裂得到的\omterm{def:g10:universal-dna:information}{染色体}已为下一轮部分复制，
分裂便可以比单次复制所需的时间
更快地相继进行。
\end{solution}

\begin{solution}{pb:g11:dna-replication:1}
\textbf{1.} 氮占质量的 16\%，并且重 7\%：
$0.16 \times 0.07 \approx 1.1\%$。梯度法可分离密度仅相差
几分之一个百分点的物质，因此 1\% 的差异
足以得到分离良好的条带。

\textbf{2.} 目的在于使几乎全部 \omterm{def:g10:universal-dna:information}{DNA} 都是重链，而不只是一半或四分之三，
即经过 14 代复制后，原始轻链所剩不到
$10^4$ 分之一。

\textbf{3.} 恰好位于两条带的正中间，其密度是两者
的平均值。

\textbf{4.} 单个细菌只含几飞克 \omterm{def:g10:universal-dna:information}{DNA}；
照片需要微克量级，因此需要数十亿个\omterm{def:g10:cells-common-unit:cell}{细胞}。

\textbf{5.} 将培养物置于重培养基中作为对照，
检查其条带始终为重带，
并在全程检查在轻培养基中生长的细菌的轻带。

\textbf{6.} 半保留复制：第 1 代，一条杂合带；
第 2 代，杂合带与轻带，数量相等。

\textbf{7.} 全保留复制：第 1 代，重带与轻带，数量相等；
第 2 代，重带与轻带之比为 1:3。

\textbf{8.} 分散复制：第 1 代，一条位于杂合位置的带；
第 2 代，一条带位于由轻到重的四分之一处。

\textbf{9.} 预测不出杂合带的全保留模型被排除；
半保留模型与分散模型都预测出单一杂合带，
二者均保留。

\textbf{10.} 半保留复制胜出：只有它预测会出现两条分离的带。
分散模型预测仅有单一条带，且逐渐变轻，
从不分裂。

\textbf{11.} 共 $2^n$ 个分子，
其中恰有 2 个含原始链。

\textbf{12.} 轻：杂合 $= (2^n - 2) : 2$：第 3 代，$3:1$；
第 4 代，$7:1$；第 6 代，$31:1$。

\textbf{13.} 杂合分子比例 $2/2^n < 0.05$，当 $2^n > 40$ 时成立：
自第 6 代起（$2/64 \approx 3\%$）。

\textbf{14.} 两条带，一重一轻，数量相等：
杂合分子由一条重链与一条轻链组成。
这一结果检验了分散模型，该模型预测每条单链都具有中间密度，
因此可独立地将其排除。

\textbf{15.} 从轻链出发，在重培养基中培养一代得到一条杂合带，
培养两代得到杂合带与重带：
与前者模式正好相反，
说明结果不取决于哪种同位素为“旧”。

\textbf{16.} $\num{4.6e6}/(2 \times 1000) = \qty{2300}{s}$，
约 38 分钟。

\textbf{17.} 77 分钟。在环状分子上，两条链在任一开口两侧均暴露；
从同一起点沿相反方向离开的两个\omterm{prop:g11:dna-replication:forks}{复制叉}，
是覆盖整个环最简单的方式。

\textbf{18.} 一个起点在 8 小时内复制 $2 \times 50 \times 8 \times 3600 =
\num{2.88e6}$ 个碱基对；至少需要 $\num{6.4e9}/\num{2.88e6}
\approx 2200$ 个起点（实际上有数万个，
它们在不同时间启动）。

\textbf{19.} 细菌：每次复制约 $\num{9.2e6} \times 10^{-9} \approx 0.01$ 个误差；
人类\omterm{def:g10:cells-common-unit:cell}{细胞}约为 6 个。人类\omterm{def:g10:universal-dna:gene}{基因}组中大部分并非\omterm{def:g10:universal-dna:gene}{基因}，
因此多数误差落在不产生影响的位置，
且\omterm{def:g10:cells-common-unit:cell}{细胞}中每个\omterm{def:g10:universal-dna:gene}{基因}都有两份拷贝。

\textbf{20.} 第 2 代的两条带——杂合带与轻带——
数量相等，证明复制为半保留方式；两个\omterm{prop:g11:dna-replication:forks}{复制叉}
约 38 分钟即可复制完细菌\omterm{def:g10:universal-dna:information}{染色体}，与细菌的世代时间大致相当，
因此复制速度决定分裂的节奏。
\end{solution}
